\subsection{Transfer to a Different Plant}
\label{sec:second-plant}
Every record so far comes from one plant family, so a method could be good at that plant without being good at power systems. We re-simulated the 12 event schedules of the held-out cohort (same sources, instants and amplitudes) with an independent plant. It has classical machines, droop governors, constant-impedance loads and the competition network, and has no excitation system. The new records pass through the same PMU operator and the same instrumental noise, so only the physics differs. The voltages at the monitored buses match the original plant to within 0.1% at rest, while the response to a switching event differs by up to a factor of two. Four normal records of 120\,s serve as controls. Every method was fixed before these records were read; learner thresholds were set on the original calibration records.

Table~\ref{tab:second-plant} gives the outcome. The static physical attributor is not affected by the change of plant: it locates \SPStaticTop{} events against \SPStaticOrig{} on the original cohort. The learners lose up to \SPDrop{} sources on the new plant and raise \SPExtrasLo{} to \SPExtrasHi{} alarms outside the scheduled events, several times more than before. Training on additional noise draws helps the phasor learner on the original cohort but not on the new plant, which shows that more noise does not replace more physics. No method locates all six events of any record, and no method raised an alarm in the normal records. The dynamic bank needs per-record banks built with the plant solver and has not been run on these records. This is the main gap of the evidence: the bank's 72 of 72 is measured on its own plant family, and a run on these records would bound how much of it depends on that match.
\begin{table}[t]\centering\scriptsize\setlength{\tabcolsep}{1.5pt}
\caption{Transfer to records simulated with a different plant (classical machines, droop governors, no excitation), built from the same 12 event schedules and passed through the same PMU operator and noise. Detected events, events located at the source (Top-1) and alarms matching no event, out of 72 events, for the original cohort and the new plant. The last column gives records with all six events located, out of 12; all methods raised zero alarms on the normal records. The dynamic bank needs per-record banks built with the original plant solver and was not run on these records. The two ``more noise'' learners were trained on four extra noise draws per training record.}
\label{tab:second-plant}
\begin{tabular}{lccccccc}\toprule
 & \multicolumn{3}{c}{Original cohort} & \multicolumn{3}{c}{Second plant} & \multicolumn{1}{c}{Records} \\
Method & Det. & Top-1 & Extras & Det. & Top-1 & Extras & (of 12) \\\midrule
Static physical & 48 & 40 & 25 & 52 & 42 & 48 & 0 \\
Trees, phasor changes & 61 & 41 & 82 & 60 & 39 & 170 & 0 \\
Trees, signatures & 58 & 49 & 30 & 68 & 45 & 156 & 0 \\
Trees, phasor changes, more noise & 63 & 48 & 94 & 66 & 42 & 273 & 0 \\
Trees, signatures, more noise & 55 & 49 & 21 & 59 & 46 & 75 & 0 \\
\bottomrule\end{tabular}\end{table}

